%----------------------------------------------------------------------
\section{Self-Enrichment and Topos Structure}\label{sec:topos}
%----------------------------------------------------------------------

The monoidal closed structure of Section \ref{sec:monoidal}\footnote{Section-level Toulmin. \emph{Move}: self-enrichment, topos structure, self-hosting. \emph{ZFC comparison}: ZFC cannot model itself without G\"{o}del coding; the framework's self-enrichment is structural (three toroid vertices). Self-hosting (the framework reconstructing its own substrate) is a fixed-point property, not a G\"{o}del trick. \emph{Residue instance}: the irremovability theorem is the residue mechanism applied at the meta-level --- removing discrimination produces a homological gap whose specified filling is discrimination.}
equips $\mathcal{C}(S,\kappa)$ with internal hom-objects.
This section shows the category is canonically self-enriched,
forms a topos with a four-valued internal logic, and
exhibits self-hosting and fixed-point properties. Detailed
verifications are in Appendix \ref{app:verifications}; here
we give each result's structural reason.

\subsection{Self-enrichment from triadic closure}

\begin{definition}[Enrichment data]\label{def:enrichment-data}
An enrichment of $\mathcal{C}$ over $(\mathcal{V}, \otimes, I)$
requires: hom-objects $[A,B] \in \mathcal{V}$, composition
morphisms $[B,C] \otimes [A,B] \to [A,C]$, and identity
morphisms $I \to [A,A]$, satisfying associativity and
unitality.
\end{definition}

\begin{theorem}[Canonical self-enrichment]\label{thm:self-enrich}\footnote{Axiom class \textsf{A} $\cdot$ depth 5 $\cdot$ \S\ref{app:self-enrichment-full}.  \emph{Warrant}: three enrichment ingredients = three vertices of the toroid. No choices required (canonical). Finite-$\kappa$: yes. Novel (three enrichment ingredients = three toroid vertices).}
$\mathcal{C}(S,\kappa)$ is canonically enriched over itself,
with hom-objects from the $\sigma$-perspective, composition
from $\tau$, and objects from $\kappa$. The three enrichment
ingredients are the three vertices of the $\GF(2)$ toroid.
\end{theorem}

\begin{proof}
Hom-objects: $[A,B] = \{d \in \kappa : \tau_d(A) = 1,\,
\sigma_d(B) = 1\}$ (Definition \ref{def:internal-hom}),
a legitimate object via the $\sigma$-perspective.
Composition: sequential filtration $(d_2, d_1) \mapsto
d_2 \circ d_1$, associative by $\wedge$. Identity: the
degenerate grade-0 element. All coherences follow from
the exterior algebra. (Full verification:
\S\ref{app:self-enrichment-full}.)
\end{proof}

\subsection{The subobject classifier}\label{sec:subobject}

\begin{definition}[Subobject classifier]\label{def:subobj-class}
$\Omega = \GF(2) \times \GF(2) = \{(0,0), (1,0), (0,1),
(1,1)\}$, the four-cell structure.
\end{definition}

\begin{theorem}[$\Omega$ is a subobject classifier]\footnote{Axiom class \textsf{A} $\cdot$ depth 2 $\cdot$ \S\ref{app:subobj-full}.  \emph{Warrant}: $\Omega = \GF(2)^2$ classifies all single-discriminator subobjects. Four-valued internal logic is Belnap FDE. Finite-$\kappa$: yes. Novel ($\Omega = \GF(2)^2$ gives four-valued logic).}
\label{thm:subobj}
For every monomorphism $m: S' \hookrightarrow A$, there
exists a unique $\chi_m: A \to \Omega$ such that $S'$ is
the pullback of $\top: 1 \to \Omega$ along $\chi_m$.
The classifying morphism sends each $a \in A$ to
$(\tau_d(a), \sigma_d(a))$.
\end{theorem}

\begin{proof}
Existence: $\chi_m$ is defined by the discriminator $d$
determining $S'$. Uniqueness: the pullback condition
determines each component. (Full verification:
\S\ref{app:subobj-full}.)
\end{proof}

The internal logic is four-valued (Belnap's FDE):
$(1,0)$ = true, $(0,1)$ = false, $(0,0)$ = unknown (gap),
$(1,1)$ = overdetermined (overlap). Boolean two-valued
logic is recovered on Boolean dimensions where every
element receives exactly one of $\tau$ or $\sigma$.

\subsection{Topos structure}

A fixed $\kappa$ may lack some finite limits. We construct
the topos as a fibered category over the poset of
$\kappa$-states.

\begin{definition}[The fibered category]\label{def:fibered}
$\mathbf{K}$ = poset of $\kappa$-states under sub-regime
inclusion. The Grothendieck construction $\int \mathcal{C}$
has objects $(\kappa, A)$ (an equivalence class at a
specific resolution) and morphisms carrying inclusions
plus morphisms in the finer category.
\end{definition}

\begin{theorem}[Fibered topos]\label{thm:topos}\footnote{Axiom class \textsf{AP} $\cdot$ depth 7 $\cdot$ \S\ref{app:power-object}, \S\ref{app:limits-exhaustive}, \S\ref{app:power-composite}, \S\ref{app:risc}.  \emph{Warrant}: finite limits at joins; $\Omega = \GF(2)^2$; power objects via monoidal closure. $\mathbf{K}$ is a modular lattice ($V^*$ is a vector space; Appendix \ref{app:risc}). [O4 resolved.] Power objects verified for all subobjects (Appendix \ref{app:power-composite}). Finite-$\kappa$: yes.}
$\int \mathcal{C}$ is an elementary topos.
\end{theorem}

\begin{proof}
\emph{Finite limits}: computed at the finite join
$\kappa_\vee$ of the stages involved. No completed infinity
needed. \emph{Subobject classifier}: $\Omega = \GF(2)^2$,
constant across fibers. \emph{Power objects}: $\Omega^B =
[B, \Omega]$ exists by monoidal closure (Theorem
\ref{thm:closure}); universal property verified in
\S\ref{app:power-object} (for single-discriminator
subobjects; composite subobjects require computing
at the join of relevant stages). The operationalist axiom
(Definition \ref{def:operationalist}) is needed only for
infinite-limit claims.
\end{proof}

\begin{corollary}[Pro-system interpretation]
\label{cor:pro-topos}
$\int \mathcal{C}$ is equivalently a pro-system of topoi
connected by geometric morphisms. The filtered colimit is
$\kappa$-indistinguishable from the pro-system
(Proposition \ref{prop:infinity}).
\end{corollary}

\subsection{Proof theory}\label{sec:proof-theory}

\begin{definition}[Evidence semantics]\label{def:evidence}
$\tau_d(a) = 1$ is a positive assertion that $\tau$ fires.
$\tau_d(a) = 0$ is silence, not falsehood. The value $0$
reports ``no match,'' not ``false.''
\end{definition}

\begin{definition}[Information ordering]\label{def:info-order}
The Belnap lattice: $(0,0) \sqsubseteq \{(1,0), (0,1)\}
\sqsubseteq (1,1)$. Gap = minimal information. Overlap =
maximal (conflicting) information. Boolean states = stable
middle.
\end{definition}

\begin{proposition}[Residues are extremal]\label{prop:extremal}
Gap $(0,0)$ and overlap $(1,1)$ are the extremes of the
information ordering. Both have $\chi = 0$ and trigger
$\kappa$-evolution. The Boolean states $(1,0)$, $(0,1)$
have $\chi = 1$ and are stable.
\end{proposition}

\begin{theorem}[Dynamic proof theory]\label{thm:proof-theory}\footnote{Axiom class \textsf{A} $\cdot$ depth 3 $\cdot$ \S\ref{app:proof-theory-full}.  \emph{Warrant}: connectives from $\GF(2)^2$ arithmetic. Excluded middle fails at the gap; explosion fails at the overlap. Finite-$\kappa$: yes. Novel (dynamic FDE with $\kappa$-evolution).}
\emph{Static}: within fixed $\kappa$, the logic is Belnap
FDE. Conjunction = componentwise multiplication. Disjunction
= componentwise maximum. Negation $\neg(\tau, \sigma) =
(\sigma, \tau)$ (swap). Excluded middle fails at the gap.
Explosion ($p \wedge \neg p \vdash q$) fails at the
overlap. Modus ponens requires Boolean premises.

\emph{Dynamic}: $\kappa$-evolution resolves residues,
monotonically increasing information. Each evolution step
moves elements from $\{(0,0), (1,1)\}$ toward
$\{(1,0), (0,1)\}$: from extremes toward the Boolean
middle.

(Full specification of inference rules and relationship
to classical, intuitionistic, and paraconsistent logic:
\S\ref{app:proof-theory-full}.)
\end{theorem}

\begin{proposition}[First-order logic is internal]\footnote{Axiom class \textsf{AP} $\cdot$ depth 8 $\cdot$ \S\ref{app:fol}.  \emph{Warrant}: quantifiers as adjunctions to pullback; requires the topos structure (class \textsf{AP}). Finite-$\kappa$: yes (via finite topos). Standard (internal logic of a topos).}
\label{prop:fol}
FOL is a fragment of the topos's internal logic:
$\forall_f$ = right adjoint to pullback $f^*$,
$\exists_f$ = left adjoint, modus ponens = evaluation
morphism, universal instantiation = $\forall$-counit,
existential generalization = $\exists$-unit. Terms are
morphisms, variables are elements, substitution is
composition. The metatheory needed to validate the
framework is available internally.
(Full construction: \S\ref{app:fol}.)
\end{proposition}

\subsection{Boolean structure and self-hosting}

\begin{proposition}[Boolean logic as a dependent type]
\label{prop:bool-dependent}
Booleanness on dimension $d$ over $S$ is the type
$\chi_d: S \to \{1\}$ (all elements discriminated). This
is a dependent type: its validity depends on $d$ and $S$,
and must be verified constructively for each dimension.
Classical logic is not a restriction on the framework but a
theorem about specific dimensions.
(Full construction: \S\ref{app:bool-convergence}.)
\end{proposition}

\begin{theorem}[Self-hosting]\label{thm:self-hosting}\footnote{Axiom class \textsf{AP} $\cdot$ depth 6 $\cdot$ \S\ref{app:self-hosting-full}.  \emph{Warrant}: Boolean restriction of $\Omega$ reconstructs $\GF(2)$ internally. \emph{Qualifier}: the self-reference is non-pathological ($\GF(2)_{\mathrm{int}}$ is algebraically identical Finite-$\kappa$: yes. Novel (Boolean restriction reconstructs $\GF(2)$ internally). to $\GF(2)_{\mathrm{ext}}$).}
The Boolean dependent type, restricted to
$\{(1,0), (0,1)\}$, reconstructs $\GF(2)$ internally:
the field with two elements, addition modulo 2,
multiplication as conjunction. The framework can serve as
its own substrate: $\GF(2)_{\mathrm{ext}} \to$
four-valued logic $\to$ Boolean type $\to
\GF(2)_{\mathrm{int}} \to \cdots$. The internal
$\GF(2)$ is algebraically identical Finite-$\kappa$: yes. Novel (Boolean restriction reconstructs $\GF(2)$ internally). to the external one.
(Proof: \S\ref{app:self-hosting-full}.)
\end{theorem}

\subsection{The fixed-point structure}

\begin{theorem}[Discrimination is a fixed point]\footnote{Axiom class \textsf{A} $\cdot$ depth 4 $\cdot$ \S\ref{app:fixed-point-full}, \S\ref{app:fixed-point-stability}.  \emph{Warrant}: three structures on the toroid; 2-cell is their binding. \emph{Qualifier}: uniqueness proved ($\dim \bigwedge^2(\GF(2)^2) = 1$). [O5 resolved.] Finite-$\kappa$: yes. Novel (toroid geometry determines the unique 2-cell).}
\label{thm:entailed}
Discrimination is a stable fixed point of the $\GF(2)$
toroid's closure operation. Three structures --- the
free-forgetful adjunction, $\GF(2)$, and Belnap logic ---
sit on the toroid as vertices. Their binding 2-cell is
discrimination. Uniqueness (that no other operation fills
the 2-cell) follows from $\dim \bigwedge^2(\GF(2)^2) = 1$: there is exactly one non-zero grade-2 element
(Appendix \ref{app:fixed-point-stability}).
(Proof: \S\ref{app:fixed-point-full}.)
\end{theorem}

\begin{theorem}[Discrimination is irremovable]\footnote{Axiom class \textsf{AP} $\cdot$ depth 8+ $\cdot$ \S\ref{app:irremovable-full}.  \emph{Warrant}: removal produces four independent failures. \emph{Qualifier}: internal necessity only --- uses the framework's own gap-filling machinery. A reader who has not accepted the framework has no reason to accept this argument. Finite-$\kappa$: yes. Novel but circular (uses framework\'s own machinery).}
\label{thm:irremovable}
Removing discrimination produces four independent failures
(sheaf, manifold, bimodule, and self-resolution) and a
homological gap whose specified filling is discrimination.
The removal reinstates what it removed. This is internal
necessity, not ontological necessity: it holds within the
framework's own machinery.
(Proof: \S\ref{app:irremovable-full}.)
\end{theorem}

\subsection{Sheaf-theoretic interpretation}\label{sec:sheaf}

\begin{definition}[The site of regimes]\label{def:site}
The site $\mathfrak{S}$ has $\kappa$-states as objects,
sub-regime inclusions as morphisms, and the Grothendieck
topology generated by the six triadic adjunctions as
covering families.
\end{definition}

\begin{theorem}[Sheaf topos]\label{thm:sheaf}\footnote{Axiom class \textsf{AP} $\cdot$ depth 6 $\cdot$ \S\ref{app:sheaf-full}.  \emph{Warrant}: six adjunctions as covering families; sheaf condition = order-independence. \emph{Qualifier}: individual adjunctions are retract adjunctions; the topology's non-triviality comes from their collective action. Finite-$\kappa$: yes. Novel (six retract adjunctions generate non-trivial topology).}
$\mathcal{C}(S,\kappa)$ at each fiber is a sheaf topos
over $\mathfrak{S}$: the triadic adjunctions provide
covering data, the sheaf condition is order-independence
(Proposition \ref{prop:monoidal}), and finite limits at
finite joins ensure compatibility.
(Proof: \S\ref{app:sheaf-full}.)
\end{theorem}

\subsection{The expanded classifier: $\GF(2)^3$}
\label{sec:expanded}

The base $\GF(2)^2$ classifier conflates
``not interrogated'' with ``interrogated, both silent.''
Adding a coverage bit $\gamma$ produces
$\GF(2)^3 = \{(\gamma, \tau, \sigma)\}$: $\gamma = 0$
means the discriminator has not been applied; $\gamma = 1$
means it has, and $(\tau, \sigma)$ records the result.

\begin{definition}[Coverage profile]\label{def:coverage}
The \emph{coverage profile} of element $a$ under
discriminator $d$ is
$(\gamma_d(a), \tau_d(a), \sigma_d(a)) \in \GF(2)^3$,
where $\gamma_d(a) = 1$ iff $d$ has been applied to $a$.
\end{definition}

\begin{theorem}[Fano closure]\label{thm:fano}\footnote{Axiom class \textsf{A} $\cdot$ depth 2 $\cdot$ \S\ref{app:fano-gluing}.  \emph{Warrant}: $PG(2, \GF(2))$ has 7 lines; triangle identity provides inference on each. Finite-$\kappa$: yes. Novel ($PG(2, \GF(2))$ from coverage bit).}
The seven non-zero points of $\GF(2)^3$ form the Fano
plane $PG(2, \GF(2))$. Each line is an instance of the
triangle identity $p_3 = p_1 + p_2$. If two points on a
Fano line have been interrogated ($\gamma = 1$), the
third is determined by inference.
(Fano line table and gluing mechanics:
\S\ref{app:fano-gluing}.)
\end{theorem}

\begin{corollary}[Self-modeling]\label{cor:self-model}
The sole true presheaf state (no inference possible) is
$(0,0,0)$: the zero element, the only point not on the
Fano plane. Every other state is either directly
interrogated or inferrable.
\end{corollary}

\begin{theorem}[Convergence rate]\label{thm:convergence}\footnote{Axiom class \textsf{A} $\cdot$ depth 3 $\cdot$ \S\ref{app:bool-convergence}.  \emph{Warrant}: 3 non-collinear points determine all 7 via Fano inference. Finite-$\kappa$: yes. Novel (three interrogations suffice).}
Three linearly independent interrogations per dimension
suffice to determine all seven Fano points: the
convergence rate is $3$.
(Proof: \S\ref{app:bool-convergence}.)
\end{theorem}
